\chapter{The \mlgw Software Library}
\label{app:ml4gw}
\begin{refsection}

\mlgw is a lightweight, PyTorch-based library designed to support the development, scaling, and deployment of machine learning (ML) models for gravitational-wave (GW) data analysis. 
Built to provide familiar functionality to anyone used to working with standard CPU-based GW libraries, \mlgw lowers the barrier to entry for GW researchers who want to incorporate ML into their work and take advantage of optimized hardware. 
Its design focuses on bridging the gap between domain-specific GW signal processing and general-purpose ML development, allowing researchers to build the tools needed to accomplish their scientific goals without getting stuck on infrastructure overhead. 
By simplifying the interface between domain science and model training, \mlgw accelerates the creation of robust, physically-informed ML pipelines for gravitational-wave astrophysics.

Machine-learning algorithms are well-suited for applications in GW astrophysics due to the field's abundance of high-quality data.
The existence of multiple independent GW detectors allows for effectively unlimited combinations of noise samples via time-shifts, and, for the most common signal morphologies, the high-fidelity simulations provided by General Relativity allow as many signal samples as desired.
However, the standard GW libraries are not designed for the scale or speed desired for ML development; conversely, standard ML libraries lack the domain-specific preprocessing functions and waveform handling tools required to train models for GW applications.

\mlgw addresses this gap by re-implementing common GW processing steps as PyTorch~\cite{paszke2019pytorch} modules, taking advantage of the natural parallelization that comes with PyTorch's batch processing, and adding the option to accelerate these steps using GPUs and other coprocessors supported by the PyTorch framework.
From libraries such as GWpy~\cite{2019ascl.soft12016M}, \mlgw re-implements power spectral density estimation, signal-to-noise ratio calculation, whitening filters, and Q-transforms.
Like bilby~\cite{Ashton_2019}, \mlgw provides the functionality to sample from astrophysical parameter distributions, which can then be used to simulate waveforms, mimicking the simulation features of lalsuite~\cite{lalsuite}.
\mlgw has available basic compact binary merger waveform families used in online searches and inference (TaylorF2, IMRPhenomD, and IMRPhenomPv2), as well as sine-Gaussian waveforms for capturing unmodeled GW signals, with more complex waveforms planned for the future.
All of these modules have been designed to work with batches of multi-channel time-series data, run on accelerated hardware, and be composable so that the output of one function can easily become the input of another.

Additionally, \mlgw contains a number of general utility features:

\begin{itemize}
    \item Efficient out-of-memory dataloading to scale the quantity of data used for training
    \item Random sampling of windows from batches of multi-channel time-series data
    \item On-the-fly signal generation for efficient scaling of training data
    \item Cubic spline interpolation to supplement PyTorch's existing interpolation functions
    \item Basic out-of-the-box neural network architectures to streamline the startup process for new users
    \item Stateful modules for handling streaming data
\end{itemize}

All implementations are fully differentiable, allowing algorithms to employ physically-motivated loss functions.
These features make it possible to train models on large quantities realistic detector data, perform data augmentation consistent with physical priors, and evaluate results in a way that is directly comparable with existing pipelines.

\mlgw is used to support the development of multiple GW analyses. 
It has been integrated into DeepClean~\cite{deepclean}, a noise-subtraction pipeline; \aframe~\cite{aframe-methods}, a search pipeline for gravitational waves from compact binary mergers; GWAK~\cite{Raikman:2023ktu}, a gravitational-wave anomaly detection pipeline; and \amplfi~\cite{amplfi}, a gravitational-wave parameter estimation pipeline.
\mlgw has enabled these algorithms to efficiently train ML models at scale and deploy models on real-time streaming data, while the use of a standardized tool set has allowed for easier communication between developers.

\singlespacing
\printbibliography[title=References, heading=subbibintoc]

\end{refsection}